% FIXME Ao final, deixe descomentada a linha correspondente ao numero de paginas
% que a sua defesa possui.
\documentclass[12pt, letterpaper, oneside]{book}  % Para menos de 100 paginas.
% \documentclass[12pt, letterpaper, twoside]{book}  % Para mais de 100 paginas.

%%%%%%%%%%%%%%%%%%%%%%%%%%%% Configuração: pacotes %%%%%%%%%%%%%%%%%%%%%%%%%%%%
%%%%%%%%%%%%%%%%%%%%%%%%%%%%%% Pacotes: básicos %%%%%%%%%%%%%%%%%%%%%%%%%%%%%% 
\usepackage[utf8]{inputenc}
\usepackage{cmap}
\usepackage[T1]{fontenc}
\usepackage[brazilian]{babel}
\usepackage{indentfirst}
\usepackage[top=3cm,bottom=3cm,right=2cm,left=2cm]{geometry}


%%%%%%%%%%%%%%%%%%%%%%%%%%%%%%% Pacotes: links %%%%%%%%%%%%%%%%%%%%%%%%%%%%%%%
\usepackage[hyphens]{url}
\usepackage{breakurl}
\usepackage{hyperref}
% FIXME Comente o pacote abaixo quando for concluir sua defesa e for entregar a
% versão final.
\usepackage{showkeys}


%%%%%%%%%%%%%%%%%%%%%%%%%%%%%%%% Pacotes: ams %%%%%%%%%%%%%%%%%%%%%%%%%%%%%%%% 
\usepackage{amsmath}
\usepackage{amsfonts}
\usepackage{amssymb}
\usepackage{amsthm}
\usepackage{breqn}


%%%%%%%%%%%%%%%%%%%%%%%%%%%%%% Pacotes: tabelas %%%%%%%%%%%%%%%%%%%%%%%%%%%%%%
\usepackage{multicol}
\usepackage{multirow}
\usepackage{array}
\usepackage{booktabs}


%%%%%%%%%%%%%%%%%%%%%%%%%%%%%% Pacotes: figuras %%%%%%%%%%%%%%%%%%%%%%%%%%%%%% 
\usepackage{pdfpages}
\usepackage{graphicx}
\usepackage{tikz}
\usepackage{wrapfig}


%%%%%%%%%%%%%%%%%%%%%%%%%%%%% Pacotes: algoritmos %%%%%%%%%%%%%%%%%%%%%%%%%%%%% 
\usepackage{algorithmic}
\usepackage[chapter]{algorithm}
\floatname{algorithm}{Algoritmo}
\renewcommand{\listalgorithmname}{Lista de Algoritmos}


%%%%%%%%%%%%%%%%%%%%%%%%%%%%%% Pacotes: códigos %%%%%%%%%%%%%%%%%%%%%%%%%%%%%% 
\usepackage{textcomp}
\usepackage{listings}
\renewcommand\lstlistingname{Código}
\renewcommand\lstlistlistingname{Lista de Códigos}


%%%%%%%%%%%%%%%%%%%%%%%%%%%%%%% Pacotes: index %%%%%%%%%%%%%%%%%%%%%%%%%%%%%%% 
\usepackage{makeidx}
\makeindex


%%%%%%%%%%%%%%%%%%%%%%%%%%%%%%% Pacotes: fontes %%%%%%%%%%%%%%%%%%%%%%%%%%%%%% 
\usepackage{lmodern}
\usepackage{mathrsfs}

%%%%%%%%%%%%%%%%%%%%%%%%%%%%%%% Pacotes: unidades %%%%%%%%%%%%%%%%%%%%%%%%%%%% 
\usepackage{siunitx}

% TODO Inserir pacotes adicionais aqui.
  % Arquivo com os pacotes.

%%%%%%%%%%%%%%%%%%%%%%%%% Configuração: dados pessoais %%%%%%%%%%%%%%%%%%%%%%%%%
% FIXME Substituir 'Nome completo do aluno' pelo seu nome.
\newcommand{\autor}{Nome completo do aluno}
% FIXME Se for do sexo feminino, descomente a linha a seguir.
% \def\femaleAuthor{}

% FIXME Substituir 'Título da defesa' pelo título da defesa.
\newcommand{\titulo}{Título da defesa}
% FIXME Se estiver no programa de mestrado, descomente a linha a seguir.
% \def\mestrado{}
% FIXME Deixe descomente apenas a linha referente ao departamento.
% \def\matematica{}
\def\aplicada{}
% \def\estatistica{}

% FIXME Substituir 'Nome completo do orientador' pelo nome completo do seu
% orientador.
\newcommand{\orientador}{Nome completo do orientador}
% FIXME Se for orientado por uma mulher, descomente a linha a seguir.
% \def\femaleOrientador{}

% FIXME Substituir 'Nome completo do coorientador' pelo nome completo do seu
% coorientador. Caso não tenha coorientador, comente a linha a seguir.
\newcommand{\coorientador}{Nome completo do coorientador}
% FIXME Se for coorientado por uma mulher, descomente a linha a seguir.
% \def\femaleCoorientador{}

% FIXME Substituir 'Ano' pelo ano em que ocorreu sua defesa.
\newcommand{\ano}{Ano}

%%%%%%%%%%%%%%%%%%%%%%%%%%% Configuração: definições %%%%%%%%%%%%%%%%%%%%%%%%%%%
%%%%%%%%%%%%%%%%%%%%%%%%%%%%% Configurações: links %%%%%%%%%%%%%%%%%%%%%%%%%%%%%
\hypersetup{
% TODO Por padrão os links, no pdf, para equações, figuras, referencias,
% tabelas, urls são identificados por uma caixa colorida em volta do link. Essa
% caixa colorida não eh impressa mas pode atrapalhar a leitura para alguns. Se
% desejar removê-las descomente a linha abaixo.
% hidelinks,
pdftitle={\titulo},  % Não modifique esta linha.
pdfauthor={\autor}  % Não modifique esta linha.
}


%%%%%%%%%%%%%%%%%%%%%%%%%% Configurações: numeração %%%%%%%%%%%%%%%%%%%%%%%%%% 
\numberwithin{equation}{section}
\numberwithin{section}{chapter}


%%%%%%%%%%%%%%%%%%%%%%%%%%%%% Configurações: ams %%%%%%%%%%%%%%%%%%%%%%%%%%%%% 
\theoremstyle{definition}
\newtheorem{thm}{Teorema}[section]
\newtheorem{con}[thm]{Conjectura}
\newtheorem{cor}[thm]{Corolário}
\newtheorem{dfn}[thm]{Definição}
\newtheorem{exm}[thm]{Exemplo}
\newtheorem{lem}[thm]{Lema}
\newtheorem{obs}[thm]{Observação}
\newtheorem{pps}[thm]{Proposição}


% TODO Inserir configurações adicionais aqui.
  % Arquivo com algumas configurações.

%%%%%%%%%%%%%%%%%%%%%%%%%% Início do texto da defesa %%%%%%%%%%%%%%%%%%%%%%%%%%%
\begin{document}
% WARNING Todas as paginas deverão ser numeradas.
%
% As paginas iniciais são numeradas com algoritmos romanos em sua forma
% minuscula.
\frontmatter
%
\thispagestyle{plain}
\includegraphics[width=.94in, height=1in,
keepaspectratio=true]{figuras/unicamp-logo.jpg}
\vspace*{1cm}
\begin{center}
  % O tamanho da fonte deve ser 16pt.
  % Deve-se utilizar caixa alta.
  {\Large \scshape \autor}
\end{center}
\vspace{4cm}
\begin{center}
  % O tamanho da fonte deve ser 16pt em negrito.
  % Deve-se utilizar caixa alta.
  {\Large \scshape \bfseries \titulo}
\end{center}
\vfill
\begin{center}
  % O tamanho da fonte deve ser 12pt em negrito.
  % Deve-se utilizar caixa alta.
  {\bfseries CAMPINAS \\ \ano}
\end{center}
  % Não edite esse arquivo.
\newpage\mbox{}\thispagestyle{plain}\newpage  % Pagina em branco.
%
% WARNING A folha de rosto precisa ser assinada pelos orientadores.
% FIXME Substitua arquivo folha-de-rosto.pdf por uma copia escaneada, comente
% esta linha e descomente a próxima.
\thispagestyle{plain}
% WARNING Não modifique este arquivo.
\includegraphics[width=.94in, height=1in,
keepaspectratio=true]{figuras/unicamp-logo.jpg}
\begin{center}
  {\large \scshape \bfseries Universidade Estadual de Campinas
  \vspace{.5cm}

  Instituto de Matemática, Estatística \\
  e Computação Científica}
\end{center}
\vspace{.7cm}
\begin{center}
  {\large \scshape \bfseries \autor}
\end{center}
\vspace{.7cm}
\begin{center}
  {\Large \scshape \bfseries \titulo}
\end{center}
\vspace{.8cm}
{\bfseries
\noindent
Orientador\ifx\femaleOrientador\undefined
\else
a\fi: Prof\ifx\femaleOrientador\undefined
\else
a\fi. Dr\ifx\femaleOrientador\undefined
\else
a\fi. \orientador
\vspace{.25cm}

\ifx\coorientador\undefined
\else
\noindent
Coorientador\ifx\femaleCoorientador\undefined
\else
a\fi: Prof\ifx\femaleCoorientador\undefined
\else
a\fi. Dr\ifx\femaleCoorientador\undefined
\else
a\fi. \coorientador
\fi
}

\vspace{.5cm}
\begin{flushright}
  \begin{minipage}[c]{.8\textwidth}
    \begin{flushright}
      \ifx\mestrado\undefined
      Tese de doutorado
      \else
      Dissertação de mestrado
      \fi
      apresentada do Instituto de \\ Matemática,
      Estatística e Computação Científica
      da Unicamp para \\ obtenção do título de
      \ifx\mestrado\undefined
      \ifx\femaleAuthor\undefined
      Doutor
      \else
      Doutora
      \fi
      \else
      \ifx\femaleAuthor\undefined
      Mestre
      \else
      Mestra
      \fi
      \fi
      em
      \ifx\matematica\undefined
      \else
      matemática.
      \fi
      \ifx\aplicada\undefined
      \else
      matemática aplicada.
      \fi
      \ifx\estatistica\undefined
      \else
      estatística.
      \fi
    \end{flushright}
  \end{minipage}
\end{flushright}
\vspace{.5cm}
\noindent
{\footnotesize \scshape
Este exemplar corresponde à versão final da 
\ifx\mestrado\undefined
tese
\else
dissertação
\fi \\
defendida 
\ifx\femaleAuthor\undefined
pelo aluno
\else
pela aluna
\fi
\autor,\\
e orientada pel\ifx\femaleOrientador\undefined
o\else
a\fi Prof\ifx\femaleOrientador\undefined
\else
a\fi. Dr\ifx\femaleOrientador\undefined
\else
a\fi. \orientador.
}
\vspace{1cm}

\noindent
{\small \bfseries
\noindent
Assinatura
\ifx\femaleOrientador\undefined
do Orientador
\else
da Orientadora
\fi

\vspace{.7cm}
\noindent
\rule[1pt]{7cm}{.5pt}  % Linha para assinatura do orientador
}
\vspace{.5cm}


\ifx\coorientador\undefined
\else
{\small \bfseries
\noindent
Assinatura
\ifx\femaleCoorientador\undefined
do Coorientador
\else
da Coorientadora
\fi

\vspace{.7cm}
\noindent
\rule[1pt]{7cm}{.5pt}  % Linha para assinatura do coorientador
}
\fi
\vfill
\begin{center}
  {\small \scshape \bfseries Campinas \\ \ano}
\end{center}

% \includepdf{folha-de-rosto.pdf}
%
% WARNING A ficha catalográfica deve estar no verso da folha de rosto.
% FIXME O arquivo ficha-catalografica.pdf deve ser sobrescrito com uma cópia
% do arquivo pdf que a biblioteca lhe enviar.
\includepdf{ficha-catalografica}
%
% WARNING A folha de aprovação deve ser assinada pelos membros da banca apos a
% defesa.
% FIXME Substitua o arquivo folha-de-aprovacao.pdf por uma copia escaneada.
\includepdf{folha-de-aprovacao}
%
% FIXME Se não for incluir a dedicatória, comentar a linha abaixo.
\chapter*{}
\begin{center}
  \emph{
  % FIXME Remover as duas linhas abaixo.
  Aos meus \ldots. \\
  Para \ldots.} 
  % TODO Inserir a dedicatória aqui.
\end{center}

%
% FIXME Se não for incluir os agradecimentos, comentar a linha abaixo.
\chapter*{Agradecimentos}
% FIXME Escrever os agradecimentos.
Aqui deve-se agradecer quem merece agradecimento

Se recebeu bolsa de algum órgão de fomento, não esqueça de agradecê-lo.

%
% FIXME Se não for incluir a epigrafe, comentar a linha abaixo.
%%TODO Escrever o modelo da epigrafe.
\chapter*{}
\vfill
\begin{flushright}
  Epígrafe \\
  Epígrafe \\
  Epígrafe \\
  Epígrafe
\end{flushright}

%
\chapter{Resumo}
% FIXME Remover as 3 linhas abaixo.
resumo resumo resumo
resumo resumo resumo
resumo resumo resumo
% WARNING O resumo deve conter no máximo 500 palavras.
% TODO Inserir o resumo em português aqui.

\vspace{.5cm}
\textbf{Palavras-chave}:
% FIXME Remover a linha abaixo.
palavra01, palavra02, palavra03.
% TODO Inserir as palavras-chave aqui.

\chapter{Abstract}
% FIXME Remover as 3 linhas abaixo.
abstract abstract abstract
abstract abstract abstract
abstract abstract abstract
% WARNING O abstract deve conter no máximo 500 palavras.
% TODO Inserir o resumo em inglês aqui.

\vspace{.5cm}
\textbf{Keywords}:
% FIXME Remover a linha abaixo.
keyword01, keyword02, keyword03.
% TODO Inserir as palavras-chave em inglês aqui.

%
% FIXME Comentar a linha abaixo se não desejar listar as figuras
% apresentadas.
\listoffigures
%
% FIXME Comentar a linha abaixo se não desejar listar as tabelas
% apresentadas.
\listoftables
%
% FIXME Comentar a linha abaixo se não desejar listar os algoritmos
% apresentados.
\listofalgorithms
%
% FIXME Comentar a linha abaixo se não desejar listar os trechos de código
% apresentados.
\lstlistoflistings
%
% FIXME Comentar a linha abaixo se não for apresentar os
% símbolos utilizados.
\chapter*{Símbolos}
\addcontentsline{toc}{chapter}{Símbolos}

\begin{description}
    % FIXME Remover os dois símbolos abaixo e incluir as que serão
    % utilizadas.
    \item[$\diamond$] Representa um diamante.
    \item[$\lhd$] Representa um deslocamento para a esquerda.
    \item[$\rhd$] Representa um deslocamento para a direita.
\end{description}

%
% FIXME Comentar a linha abaixo de não for apresentar as
% abreviações utilizadas.
\chapter*{Abreviações}
\addcontentsline{toc}{chapter}{Abreviações}

\begin{description}
  % FIXME Remover as duas abreviações abaixo e incluir as que serão
  % utilizadas.
  \item[FIXME] Indica que algo deve ser consertado.
  \item[TODO] Indica que algo deve ser feito.
\end{description}

%
\tableofcontents
%
% As paginas com o conteúdo da tese são numeradas com algoritmos arábicos.
\mainmatter
%
% FIXME Remover as 3 linhas abaixo.
\input{ex_cap1}
\input{ex_cap2}
\input{ex_cap3}
\input{ex_cap4}
% TODO Inserir os arquivos referentes ao corpo da tese.
%
%
% FIXME Comentar a linha abaixo se não for apresentar o glossário.
\chapter{Glossário}

\begin{description}
  % FIXME Remover as duas explicações abaixo e incluir as que serão
  % utilizadas.
  \item[Tese] Proposição que é enunciada e sustentada.
  \item[Dissertação] Discurso, exposição ou exame minucioso de
    determinado assunto.
\end{description}

%
% FIXME Se não for utilizar apêndices, comentar a linha abaixo.
\appendix
% FIXME Remover a linha abaixo.
\input{ex_ape1}
\input{ex_ape2}
\chapter{Uso de Licenças Creative Commons}
As licenças Creative Commons são uma família de licenças copyleft criada e
mantida pela organização não comercial Creative Commons. Atualmente as licenças
Creative Commons Atribuição 3.0 Não Adaptada (CC-BY) e Creative Commons
Atribuição-CompartilhaIgual 3.0 Não Adaptada (CC-BY-SA) são as licenças mais
utilizadas em conteúdos/recursos abertos.

Os mantenedores deste modelo sugerem a utilizado das licenças CC-BY e CC-BY-SA
por acreditarem que a ciência é melhor desenvolvida quando:
\begin{itemize}
  \item pesquisadores possam compartilhar, utilizar e remixar o trabalho de seus
    colegas (sob a condição de dar os devidos créditos) e
  \item não tenham que se preocupar com a violação de direitos autorais (desde
    que faça atribuição ao autor e trabalho original).
\end{itemize}
Essas são as características de materiais licenciados utilizando CC-BY ou
CC-BY-SA e por isso incentivamos seu uso.

Nas próximas seções você encontrará maiores detalhes sobre as licenças CC-BY e
CC-BY-SA, como utilizá-las em sua dissertação de mestrado ou tese de doutorado e
informações adicionais.

\section{Creative Commons Atribuição 3.0 Não Adaptada}
A licença Creative Commons Atribuição 3.0 Não Adaptada (CC-BY) é a menos restritiva das
licenças Creative Commons. Um material que utiliza essa licença pode ser:
\begin{itemize}
  \item copiado,
  \item distribuído e
  \item modificado
\end{itemize}
por qualquer pessoa desde que indique a obra original e seu autor. Para maiores
informações visite
\url{http://creativecommons.org/licenses/by/3.0/deed.pt_BR}.


Para utilizá-la na sua dissertação de mestrado ou tese de doutorado,
certifique-se do arquivo \lstinline+cc-by.tex+ ter sido incluso no arquivo
\lstinline+tese.tex+ como indicado abaixo:
\begin{lstlisting}
\annex
% TODO Inserir os arquivos referentes aos apendices e anexos.
\chapter{Licença}
Copyright (c) \ano \ de \autor.

Esta obra está licenciada sob a licença Creative Commons Atribuição 3.0 Não
Adaptada. Para ver uma cópia desta licença, visite
\url{http://creativecommons.org/licenses/by/3.0/}.

\begin{center}
  \includegraphics{figuras/cc-by.png}
\end{center}

\section{Sobre a licença dessa obra}
A licença Creative Commons Atribuição 3.0 Não Adaptada utilizada nessa obra diz
que:
\begin{enumerate}
  \item Você tem a liberdade de:
    \begin{itemize}
      \item Compartilhar — copiar, distribuir e transmitir a obra;
      \item Remixar — criar obras derivadas;
      \item fazer uso comercial da obra.
    \end{itemize}
  \item Sob as seguintes condições:
    \begin{itemize}
      \item Atribuição — Você deve creditar a obra da forma especificada pelo
        autor ou licenciante (mas não de maneira que sugira que estes concedem
        qualquer aval a você ou ao seu uso da obra).
    \end{itemize}
\end{enumerate}

%
\backmatter
\end{lstlisting}

\section{Creative Commons Atribuição-CompartilhaIgual 3.0 Não Adaptada}
A licença Creative Commons Atribuição-CompartilhaIgual 3.0 Não Adaptada
(CC-BY-SA) consiste na licença CC-BY acrescida de uma restrição: uma obra que
tiver sido produzida com base em uma outra obra licenciada com CC-BY-SA precisa
ser necessariamente licenciada com CC-BY-SA. Para maiores informações visite
\url{http://creativecommons.org/licenses/by-sa/3.0/deed.pt_BR}.

Para utilizá-la na sua dissertação de mestrado ou tese de doutorado,
certifique-se do arquivo \lstinline+cc-by-sa.tex+ ter sido incluso no arquivo
\lstinline+tese.tex+ como indicado abaixo:
\begin{lstlisting}
\annex
% TODO Inserir os arquivos referentes aos apendices e anexos.
\chapter{Licença}
Copyright (c) \ano \ de \autor.

Esta obra está licenciada sob a licença Creative Commons
Atribuição-CompartilhaIgual 3.0 Não Adaptada. Para ver uma cópia desta licença,
visite \url{http://creativecommons.org/licenses/by-sa/3.0/}.

\begin{center}
  \includegraphics{figuras/cc-by-sa.png}
\end{center}

\section{Sobre a licença dessa obra}
A licença Creative Commons Atribuição-CompartilhaIgual 3.0 Não Adaptada
utilizada nessa obra diz que:
\begin{enumerate}
  \item Você tem a liberdade de:
    \begin{itemize}
      \item Compartilhar — copiar, distribuir e transmitir a obra;
      \item Remixar — criar obras derivadas;
      \item fazer uso comercial da obra.
    \end{itemize}
  \item Sob as seguintes condições:
    \begin{itemize}
      \item Atribuição — Você deve creditar a obra da forma especificada pelo
        autor ou licenciante (mas não de maneira que sugira que estes concedem
        qualquer aval a você ou ao seu uso da obra).
      \item Compartilhamento pela mesma licença — Se você alterar, transformar
        ou criar em cima desta obra, você poderá distribuir a obra resultante
        apenas sob a mesma licença, ou sob uma licença similar à presente.
    \end{itemize}
\end{enumerate}

%
\backmatter
\end{lstlisting}


\section{Informações adicionais}
Perguntei para a Agência de Inovação da UNICAMP sobre quem detém os direitos
autorais das dissertações de mestrado e teses de doutorado e recebi a seguinte
resposta:

\begin{verbatim}
Date: Fri, 8 Mar 2013 08:41:22 -0300
From: Gabriel Gustavo Guion <gguion@inova.unicamp.br>
To: Raniere Silva <ra092767@ime.unicamp.br>
Cc: patentes@inova.unicamp.br
Subject: Re: Direito Autoral de Dissertações/Teses

Prezado Raniere,

Todas as obras, literárias, artísticas ou científicas são passiveis de
serem protegidas por Direito Autoral. Portanto, sua tese/dissertação é sim
registrável.
Para maiores informações acesse:
http://www.bn.br/portal/?nu_pagina=28

No caso da Unicamp, segundo nossa política de Propriedade Intelectual
instaurada pela deliberação *CONSU*-A-016/2010 de 30/11/2010, os direitos
de propriedade das criações literárias, artísticas e pedagógicas
pertencerão aos autores.

Estamos a disposição.

Atenciosamente.
Gabriel Guion
Gerente de Propriedade Intelectual

Agência de Inovação Inova Unicamp
(19) 3521-2549
www.inova.unicamp.br
\end{verbatim}

Deste modo, você como autor da sua dissertação de mestrado ou tese de
doutorado pode optar por utilizar uma licença copyleft, por exemplo CC-BY ou
CC-BY-SA, para a mesma.

É comum as informações de licenciamento de materiais impressos vir logo no
início do mesmo como é o caso na maioria dos livros (muitas vezes na mesma
página contendo a ficha catalográfica). Infelizmente como atualmente não existe
nenhuma normativa da UNICAMP especificando a posição da menção dos direitos
autorais das dissertações de mestrado e teses de doutorado acreditamos que o
melhor lugar para conter essa informação seja o último anexo\footnote{Isso
baseado no uso da licença GNU Free Documentation License (para maiores
informações visite \url{http://www.gnu.org/copyleft/fdl.html}}.

% TODO Inserir os arquivos referentes aos apêndices.
%
% TODO Se não for utilizar anexos nem uma licença Creative Commons, comentar a linha abaixo.
\annex
% FIXME Remover a linha abaixo.
\input{ex_ane1}
% TODO Inserir os arquivos referentes aos anexos.
%
% FIXME Escolher, opcionalmente, uma das licença Creative Commons abaixo.
% Para isso remova a linha correspondente a licença não desejada.
\chapter{Licença}
Copyright (c) \ano \ de \autor.

Esta obra está licenciada sob a licença Creative Commons Atribuição 3.0 Não
Adaptada. Para ver uma cópia desta licença, visite
\url{http://creativecommons.org/licenses/by/3.0/}.

\begin{center}
  \includegraphics{figuras/cc-by.png}
\end{center}

\section{Sobre a licença dessa obra}
A licença Creative Commons Atribuição 3.0 Não Adaptada utilizada nessa obra diz
que:
\begin{enumerate}
  \item Você tem a liberdade de:
    \begin{itemize}
      \item Compartilhar — copiar, distribuir e transmitir a obra;
      \item Remixar — criar obras derivadas;
      \item fazer uso comercial da obra.
    \end{itemize}
  \item Sob as seguintes condições:
    \begin{itemize}
      \item Atribuição — Você deve creditar a obra da forma especificada pelo
        autor ou licenciante (mas não de maneira que sugira que estes concedem
        qualquer aval a você ou ao seu uso da obra).
    \end{itemize}
\end{enumerate}
  % Creative Commons Atribuição 3.0 Não Adaptada
\chapter{Licença}
Copyright (c) \ano \ de \autor.

Esta obra está licenciada sob a licença Creative Commons
Atribuição-CompartilhaIgual 3.0 Não Adaptada. Para ver uma cópia desta licença,
visite \url{http://creativecommons.org/licenses/by-sa/3.0/}.

\begin{center}
  \includegraphics{figuras/cc-by-sa.png}
\end{center}

\section{Sobre a licença dessa obra}
A licença Creative Commons Atribuição-CompartilhaIgual 3.0 Não Adaptada
utilizada nessa obra diz que:
\begin{enumerate}
  \item Você tem a liberdade de:
    \begin{itemize}
      \item Compartilhar — copiar, distribuir e transmitir a obra;
      \item Remixar — criar obras derivadas;
      \item fazer uso comercial da obra.
    \end{itemize}
  \item Sob as seguintes condições:
    \begin{itemize}
      \item Atribuição — Você deve creditar a obra da forma especificada pelo
        autor ou licenciante (mas não de maneira que sugira que estes concedem
        qualquer aval a você ou ao seu uso da obra).
      \item Compartilhamento pela mesma licença — Se você alterar, transformar
        ou criar em cima desta obra, você poderá distribuir a obra resultante
        apenas sob a mesma licença, ou sob uma licença similar à presente.
    \end{itemize}
\end{enumerate}
  % Creative Commons AtribuiçãoCompartilhaIgual 3.0 Não Adaptada
%
\backmatter
\printbibliography
\addcontentsline{toc}{chapter}{Referência Bibliográfica}
%
\clearpage
\phantomsection
\addcontentsline{toc}{chapter}{Índice Remissivo}
\printindex
\end{document}
